\begin{helper}
For any event \(A\) and conditioning event \(C\), suppose all
\(\noderef{TwoBiteSampleWeight}\) values are nonnegative.  If
\[
\noderef{TwoBiteEventProbability}(C)\le B_C,\qquad
\noderef{TwoBiteConditionalProbability}(A\mid C)\le B_A,
\]
with \(B_C\ge0\) and \(B_A\ge0\), then
\[
\noderef{TwoBiteEventProbability}(A\cap C)\le
B_C\min(1,B_A).
\]
\end{helper}
\begin{proof}
Write \(P(E)\) for \(\noderef{TwoBiteEventProbability}(E)\).  Because all
sample weights are nonnegative, \(P(E)\ge0\) for every event \(E\), and
monotonicity holds: if \(E\subseteq F\), then \(P(E)\le P(F)\).  Hence
\(P(A\cap C)\le P(C)\), giving the factor \(1\).

If \(P(C)=0\), monotonicity and nonnegativity give \(P(A\cap C)=0\), so the
claim follows.  If \(P(C)>0\), unfolding
\(\noderef{TwoBiteConditionalProbability}\) gives
\[
P(A\cap C)/P(C)\le B_A,
\]
and multiplying by the positive denominator yields
\[
P(A\cap C)\le P(C)B_A.
\]
Combining this with \(P(A\cap C)\le P(C)\) gives
\[
P(A\cap C)\le P(C)\min(1,B_A).
\]
Finally \(P(C)\le B_C\) and \(\min(1,B_A)\ge0\), so multiplication preserves
the inequality and gives the stated bound.
\end{proof}
